\documentclass[11pt,a4paper]{article}
\usepackage[margin=2.2cm]{geometry}
\usepackage{microtype}
\setlength{\parindent}{0pt}
\setlength{\parskip}{6pt}

\begin{document}
\begin{center}
{\Large\bfseries PhotonicsAI --- AI-Driven Science and Autonomous Laboratories}\\[3pt]
{\large Homework 3: Hardware in the Loop}\\[6pt]
Name: Erik Mirzoyan
\end{center}

\section*{The device}
A lab camera: the \textbf{Hikrobot MV-CS028-10UM} industrial camera (monochrome, USB~3.0, USB3~Vision), controlled
through Hikrobot's MVS software development kit (SDK).

\section*{Project description}
Until now the agent system from Homeworks~1 and~2 has worked only with code and with a simulated optical bench.
Connecting the camera gives it a real instrument. The camera's controls (exposure, gain, region of interest, frame
capture and image measurements) become tools on an MCP server, and one agent, the executor, holds them. A new
hardware turn in the graph closes the loop: the agent decides what to do, commands the camera, reads the measurement back,
and decides again. The human stop from Homework~2 stays for setting changes and for actions only a person can do, such as
blocking the light.

With this connection the system will be able to carry out a measurement task from a plain-language request, choosing the
settings itself and checking the result on the real image. It will set the exposure automatically so that the image uses
most of the camera's range without saturating; characterise the camera's noise (read noise, conversion gain, full-well
capacity) from dark frames and a photon-transfer series; find an exposure free of the room lights' flicker (mains lamps
flicker at 100\,Hz); and check the stability of the intensity over time. The lab has no laser for this homework, so the
camera looks at a white surface under the room lights. The measured camera parameters will then be put into the camera
model of the Homework~2 simulated bench, so that the agent predicts the focusing enhancement a wavefront-shaping experiment
would reach with this camera: a first step towards running the Homework~2 focusing experiment on real hardware, with the
lab's DMDs.
\end{document}
